% \iffalse meta-comment
%<*preamble>
%% SPDX-FileCopyrightText: 2026 Ruiyuan Chen
%% SPDX-License-Identifier: MIT
%%
%% ---------------------------------------------------------------
%% luaalignstate --- Access the alignment counter in LuaTeX
%% ---------------------------------------------------------------
%%
%</preamble>
%<*internal>
\expandafter\ifx\csname LaTeXReleaseInfo\endcsname\relax
    \def\endinstall{\endbatchfile}
\else
    \begingroup\let\endinstall\endgroup
\fi
%</internal>
%<*install>
\input docstrip.tex
\keepsilent
\askforoverwritefalse
\def\MetaPrefix{\noexpand\CurrentMetaPrefix}
\preamble
\endpreamble
\usepreamble\empty
\usepostamble\defaultpreamble
\usedir{tex/luatex/\jobname}
\let\CurrentMetaPrefix\DoubleperCent
\generate{\file{\jobname.sty}{
    \from{\jobname.dtx}{preamble,package}
}}
%</install>
%<*internal>
\usedir{source/luatex/\jobname}
\generate{\file{\jobname.ins}{
    \from{\jobname.dtx}{preamble,install}
}}
%</internal>
%<*install>
\def\CurrentMetaPrefix{--}
\generate{\file{\jobname.lua}{
    \from{\jobname.dtx}{preamble,lua}
}}
%</install>
%<install>\endbatchfile
%<*internal>
\endinstall
%</internal>
%<*driver>
\documentclass{l3doc}
\usepackage{microtype}
\usepackage[l3doc]{fancyvrbx}
\hypersetup{bookmarksopenlevel=2}
\DisableCrossrefs
\ExplSyntaxOn
\makeatletter
\sys_if_engine_luatex:T {
    \usepackage{\jobname}
}
\bool_set_false:N \l__codedoc_allow_indexing_bool
\cs_gset_protected:Npn \@wrindex #1 {}
\cs_gset_protected:Npn \codeline@wrindex #1 {}
\cs_gset_protected:Npn \__codedoc_index_page_hc:nn #1 #2 {}
\cs_gset_protected:Npn \__codedoc_index_codeline_hc:nn #1 #2 {}
\NewDocumentCommand \metacode {m} {
    \textnormal {\textlangle \emph{#1} \textrangle}
}
\makeatother
\ExplSyntaxOff
\begin{document}
\DocInput{\jobname.dtx}
\end{document}
%</driver>
% \fi
%
% \title{\pkg{luaalignstate} --- Access the alignment counter in \LuaTeX{}}
% \author{Ruiyuan Chen}
% \date{}
% \maketitle
%
%
% In \TeX{}, the alignment brace counter, also known as ``master counter'' (according to the \TeX{}book) or |align_state| (according to the source code), is an internal variable that is used to track the balance between all braces |{| and |}| seen so far, and plays a subtle role in the handling of |\halign|/|\valign| and derived macros such as \LaTeX{}'s |tabular|.
% The variable is not exposed to the macro language or even to the Lua libraries in \LuaTeX{}.
% This makes certain rare bugs difficult to track down,\footnote{See e.g., \url{https://github.com/latex3/latex2e/issues/1468}.}
% and also prevents some useful alignment tricks that would otherwise be possible.\footnote{For example, it is difficult to create a fully general mechanism that allows a colspan to be specified in a table cell other than at the very beginning.}
%
% This package provides a hack in \LuaTeX{} to directly inspect and modify the alignment brace counter.
% There is an underlying format-independent Lua module, as well as a \LaTeX3 programming wrapper interface.
%
%
% \section{Lua interface}
%
% The Lua module is loaded with
% \begin{verbatim}
% local luaalignstate = require"luaalignstate"
% \end{verbatim}
% which returns a table containing the following functions.
%
% \begin{function}{luaalignstate.getset()}
% \begin{latexcodesyntax}
% new, old = luaalignstate.getset(function (old_alignstate)
%     return new_alignstate
% end)
% \end{latexcodesyntax}
% The main interface.
% Determines the current alignment counter, and optionally updates it in one go.
% If the provided function is or returns |nil|, no update is made.
%
% \begin{texnote}
% This function works by executing an empty |\halign|.
% As such, it \emph{must not be called} while scanning the preamble of another alignment (via e.g., |\span\directlua|); doing so would lead to memory corruption.
%
% Because of this, the function will throw a fatal error if it detects that it was likely called from an alignment preamble, namely if |align_state| $< -500000$.
% (This will also happen outside of alignment preambles, if you for some reason set |align_state| to a value less than $-500000$.)
% \end{texnote}
% \end{function}
%
% \begin{function}{
%     luaalignstate.get(),
%     luaalignstate.set(),
%     luaalignstate.add(),
% }
% \begin{verbcodesyntax}
% alignstate = luaalignstate.get()
% new, old = luaalignstate.set(new_alignstate)
% new, old = luaalignstate.add(increment)
% \end{verbcodesyntax}
% Simple wrappers around |getset()|.
% \end{function}
%
%
% \section{\LaTeX3 interface}
%
% The \LaTeX{} package is loaded with
% \begin{expl3code}
% \RequirePackage{luaalignstate}
% \end{expl3code}
% This loads the Lua module globally, and defines the following L3 programming functions.
%
% \begin{function}[EXP]{\luaalignstate_value:}
% Expands in one step to the current value of the alignment counter (in decimal).
%
% \begin{texnote}
% As with the Lua functions, this is \emph{not} safe to expand (via |\span|) while scanning the preamble of an alignment.
% \end{texnote}
% \end{function}
%
% Note that \cs{luaalignstate_value:} can be somewhat tricky to expand correctly in the desired ambient brace level:
% for example,
% \begin{expl3code}
%     \tl_set:No \l_my_tl {\luaalignstate_value:}
% \end{expl3code}
% would likely set \cs{l_my_tl} to one or two more than the alignment counter value at the point that it is executed, depending on the implementation of \cs{tl_set:No}.
%
% For this reason, the following getter functions are provided as an alternative.
%
% \begin{function}{\luaalignstate_get:N}
% \begin{expl3codesyntax}
% \luaalignstate_get:N \meta{tl var}
% \end{expl3codesyntax}
% Sets the \meta{tl var} (locally) to the current value of the alignment counter.
% \end{function}
%
% \begin{function}[EXP]{\luaalignstate_use:n}
% \begin{expl3codesyntax}
% \luaalignstate_use:n \Arg{code}
% \end{expl3codesyntax}
% Expands in two steps to \meta{code} followed by a brace group containing the current value of the alignment counter.
% \end{function}
%
% \begin{function}[EXP]{\luaalignstate_use_inline:n}
% \begin{expl3codesyntax}
% \luaalignstate_use_inline:n \Arg{inline function}
% \end{expl3codesyntax}
% Expands in two steps to \meta{inline function} applied to the alignment counter value.
% \end{function}
%
% \begin{function}{\luaalignstate_set:n}
% \begin{expl3codesyntax}
% \luaalignstate_set:n \Arg{int expr}
% \end{expl3codesyntax}
% Sets the alignment counter to the given \meta{int expr}.
%
% \begin{texnote}
% The \meta{int expr} may refer to the previous alignment counter value via \cs{luaalignstate_value:}, which within the \meta{int expr} expands to the value at the point where \cs{luaalignstate_set:n} is called (i.e., outside any braces surrounding or part of \meta{int expr}).
%
% On the other hand, the \meta{int value} should not refer directly to \cs{currentgrouplevel}, \cs{currentgrouptype}, or \cs{lastnodetype}, as their values will be wrong.
% \end{texnote}
% \end{function}
%
% \begin{function}{\luaalignstate_add:n}
% \begin{expl3codesyntax}
% \luaalignstate_add:n \Arg{int expr}
% \end{expl3codesyntax}
% Adds \meta{int expr} to the alignment counter.
%
% \begin{texnote}
% This may be expressed in terms of \cs{luaalignstate_set:n}, but is $\approx 2$\,ops faster.
% Unlike above, here the \meta{int expr} should \emph{not} refer to \cs{luaalignstate_value:}.
% \end{texnote}
% \end{function}
%
%
% \begin{implementation}
% \newpage
% \section{Implementation}
%
%    \begin{macrocode}
%<*lua>
%<@@=luaalignstate>
%    \end{macrocode}
%
%
% \subsection{Global imports}
%
% Import needed global functions.
%    \begin{macrocode}
local pairs = pairs
local select = select
%    \end{macrocode}
% Import \LuaTeX{} modules.
%    \begin{macrocode}
local luatex = tex
local luatex_token = token
local luatex_node = node
local luatex_status = status
_ENV = nil
%    \end{macrocode}
% Import commonly used functions from modules.
%    \begin{macrocode}
local tex_getnest = luatex.getnest
local tex_runtoks = luatex.runtoks
local tex_quittoks = luatex.quittoks
local token_command_id = luatex_token.command_id
local token_new = luatex_token.new
local token_get_next = luatex_token.get_next
local token_scan_code = luatex_token.scan_code
local token_put_next = luatex_token.put_next
local token_unchecked_put_next = luatex_token.unchecked_put_next
local node_direct_new = luatex_node.direct.new
local node_direct_free = luatex_node.direct.free
local node_direct_getnext = luatex_node.direct.getnext
local node_direct_getdata = luatex_node.direct.getdata
local node_direct_setnext = luatex_node.direct.setnext
local node_direct_setdata = luatex_node.direct.setdata
local node_direct_setfield = luatex_node.direct.setfield
%    \end{macrocode}
%
%
% \subsection{Constants}
%
% \begin{macro}[int]{EQ_...}
% Some hardcoded eqtb locations (from \file{equivalents.h}).
%    \begin{macrocode}
local EQ_frozen_control_sequence = ((1 +1)+luatex_status.hash_size)
local EQ_local_base = ((((((EQ_frozen_control_sequence+13 )+(5535-0 +1))+1)+19)+65536)+65536)
%    \end{macrocode}
% \end{macro}
%
% \begin{macro}[int]{TOKEN_...}
% Commonly used tokens.
%    \begin{macrocode}
local CS_TOKEN_FLAG = 0x1FFFFFFF
local COMMAND_lua_local_call = token_command_id"lua_local_call"
local TOKEN_bgroup = token_new(("{"):byte(), 1)
local TOKEN_egroup = token_new(("}"):byte(), 2)
local TOKEN_param = token_new(("#"):byte(), 6)
local TOKEN_immediateassignment = token_new(5, token_command_id"convert")
local TOKEN_nonstopmode = token_new(1, token_command_id"set_interaction")
local TOKEN_halign = token_new(0, token_command_id"halign")
local TOKEN_span = token_new(0x110002, 4)
%    \end{macrocode}
% From \file{equivalents.h}.
% (NOTE: |\cr| must be a named control sequence, not a nameless internal token, since the internal token is treated as |out_param| when read.)
%    \begin{macrocode}
local TOKEN_cr = token_new(CS_TOKEN_FLAG + EQ_frozen_control_sequence + 1, 0)
local TOKEN_endv = token_new(CS_TOKEN_FLAG + EQ_frozen_control_sequence + 6, 0)
local TOKEN_everycr = token_new(EQ_local_base + 8, token_command_id"assign_toks")
%    \end{macrocode}
% \end{macro}
%
% \begin{macro}[int]{MODE_hmode, MODE_vmode}
% Horizontal and vertical mode constants (in |tex.nest.mode|).
%    \begin{macrocode}
local MODE_vmode
for k, v in pairs(luatex.getmodevalues()) do
    if v == "vertical" then MODE_vmode = k end
end
%    \end{macrocode}
% \end{macro}
%
% \begin{macro}[int]{NODEID_...}
% Node types.
%    \begin{macrocode}
local NODEID_align_stack = luatex_node.id"align_stack"
%    \end{macrocode}
% \end{macro}
%
%
% \subsection{Helpers}
%
% \begin{macro}[int]{tmp_whatsit_ptr}
% A whatsit to help with allocations.
%    \begin{macrocode}
local tmp_whatsit_ptr = node_direct_new("whatsit", "user_defined")
%    \end{macrocode}
% \end{macro}
%
% \begin{macro}[int]{new_lualocal}
% Creates a new (expandable) luacall that's stored in the Lua C registry.
%    \begin{macrocode}
local function new_lualocal(f)
%    \end{macrocode}
% It is necessary to set the whatsit integer value to 0 first,
% or else setting the lua value will kill any existing registry entry.
%    \begin{macrocode}
    node_direct_setfield(tmp_whatsit_ptr, "type", 100)  -- ("d"):byte()
    node_direct_setdata(tmp_whatsit_ptr, 0)
    node_direct_setfield(tmp_whatsit_ptr, "type", 108)  -- ("l"):byte()
    node_direct_setdata(tmp_whatsit_ptr, f)
    node_direct_setfield(tmp_whatsit_ptr, "type", 100)  -- ("d"):byte()
    return token_new(node_direct_getdata(tmp_whatsit_ptr), COMMAND_lua_local_call)
end
%    \end{macrocode}
% \end{macro}
%
% \begin{macro}[int]{fatal_error()}
% Raises a fatal error.
% We could use any of the several tokens that triggers a |normal_error()|, |fatal_error()|, or |confusion()|.
% Since this is called when the |align_state| is manipulated from alignment preambles, it seems that the ``interwoven alignment preamble'' error message is the most appropriate.
% The |scan_code()| is used to flush the two token lists inserted by |runtoks()| from the error context.
% (Unfortunately, there seems to be no way to also flush the pushed |\endtemplate|.)
%    \begin{macrocode}
local fatal_error do
local fatal_error_msg
local fatal_error_help
local function fatal_error_run()
    token_get_next()
    token_put_next(TOKEN_endv)
    if luatex.interactionmode > 1 then
        token_put_next(TOKEN_immediateassignment, TOKEN_nonstopmode)
    end
    token_scan_code(0)
    return luatex.error(fatal_error_msg, fatal_error_help)
end
function fatal_error(msg, help)
    fatal_error_msg = msg
    fatal_error_help = help or {}
    return tex_runtoks(fatal_error_run)
end
end
%    \end{macrocode}
% \end{macro}
%
%
% \subsection{Getting and setting the alignment counter}
%
% The module table.
%    \begin{macrocode}
local luaalignstate = {}
%    \end{macrocode}
%
% Now for the main hack.
% The strategy is to execute (the equivalent of)
% \begin{expl3code}
%     \{\metacode{determine \texttt{align_stack} location}}
%     \halign{%
%         \span\immediateassignment\everycr{}%
%         #\cr
%         \{\metacode{read/write \texttt{align_state} from \texttt{align_stack}}}
%     }
% \end{expl3code}
% When \LuaTeX{} executes |\halign|, it creates a node of type |align_stack| (see \file{align.c}:|push_alignment()|), which contains a field that records the |align_state| before the alignment began, which is restored after the alignment ends.
% Thus, the goal of the expandable luacall
% \metacode{read/write \texttt{align_state} from \texttt{align_stack}}
% is to access this field.
% Note that because it is expandable, and we nullified the |\everycr| (\emph{inside} the outer |\halign| group, but outside of any cell), the alignment won't have any rows: everything takes place while |align_peek()| is looking for the first row, only to find the end-brace |}|.
%
% Unfortunately, the |align_stack| does not seem to be referenced from any nodes in the current list, so we cannot directly access it.
% Instead we take advantage of the way that node allocation works in \LuaTeX{}: right before the |\halign|, we allocate a |align_stack| ourselves, record its address, and then free it, putting it onto the free list and ensuring that |\halign| will allocate in the same location.
% (So this is an intentional use-after-free!)
%
% We export a main public function, which needs to pass arguments to several auxiliaries via upvariables;
% we isolate these inside a |do| block for clarity.
%    \begin{macrocode}
local getset do
%    \end{macrocode}
%
% \begin{macro}[int]{getset_align_stack_ptr}
% The location of the |align_stack| node (an integer/|node.direct| pointer), determined right before the |\halign| as described above.
%    \begin{macrocode}
local getset_align_stack_ptr
%    \end{macrocode}
% \end{macro}
%
% \begin{macro}[int]{getset_func, getset_oldval, getset_newval}
% Used to pass the user-provided getset function and the old and new |align_state| values.
%    \begin{macrocode}
local getset_func, getset_oldval, getset_newval
%    \end{macrocode}
% \end{macro}
%
% \begin{macro}[int]{getset_do}
% Executed expandably inside of the |\halign|, right before it ends (without any rows).
% Fetch the enclosing |align_state| from the |align_stack| node,
% run the caller-provided update function,
% and set the |align_state| to the new value if different.
%    \begin{macrocode}
local function getset_do()
%    \end{macrocode}
% The following calls |vlink(getset_align_stack_ptr + 3)|, which stores the enclosing |align_state|;
% see
% \file{lnodelib.c}:|lua_nodelib_direct_getnext()|
% and
% \file{align.c}:|push_alignment()|.
%    \begin{macrocode}
    local oldval = node_direct_getnext(getset_align_stack_ptr + 3) or 0
    if oldval < -500000 then
        return fatal_error(
            "luaalignstate Fatal Error: Cannot access from alignment preamble",
            {
                "Accessing the alignment counter from a preamble causes memory corruption.",
                "I'm going to quit now.",
            }
        )
    end
    getset_oldval = oldval
    local newval = getset_func and getset_func(oldval)
    if newval == nil then
        getset_newval = oldval
    else
        if newval ~= oldval then
            node_direct_setnext(getset_align_stack_ptr + 3, newval)
        end
        getset_newval = newval
    end
    token_put_next(TOKEN_egroup)
    token_scan_code(0)
    return tex_quittoks()
end
local TOKEN_getset_do = new_lualocal(getset_do)
%    \end{macrocode}
% \end{macro}
%
% \begin{macro}[int]{TOKENS_getset}
% The token list containing |\halign|, sans the close brace which is inserted by |getset_do()|.
%    \begin{macrocode}
local TOKENS_getset = {
    TOKEN_halign, TOKEN_bgroup,
        TOKEN_span, TOKEN_immediateassignment, TOKEN_everycr, TOKEN_bgroup, TOKEN_egroup,
        TOKEN_param, TOKEN_cr,
        TOKEN_getset_do,
}
%    \end{macrocode}
% \end{macro}
%
% \begin{macro}[int]{getset_run}
% The function that's executed by |tex.runtoks()|.
%    \begin{macrocode}
local function getset_run()
%    \end{macrocode}
% Remove the |\endlocalcontrol|.
% This allows the user function to scan the token stream (see e.g., \cs{@@_set:w}),
% which is also why we left off the close brace above.
%    \begin{macrocode}
    token_get_next()
%    \end{macrocode}
% Set the mode to internal vertical, to prevent |\halign| from complaining.
% The enclosing mode is already remembered and restored by |runtoks()|.
%    \begin{macrocode}
    tex_getnest(tex_getnest("ptr")).mode = -MODE_vmode
%    \end{macrocode}
% Now put an |align_stack| on the free list, record its address, and insert the |halign| tokens.
%    \begin{macrocode}
    getset_align_stack_ptr = node_direct_new(NODEID_align_stack)
    node_direct_free(getset_align_stack_ptr)
    token_unchecked_put_next(TOKENS_getset)
end
%    \end{macrocode}
% \end{macro}
%
% \begin{macro}{
%     luaalignstate.getset(),
% }
% The actual main function.
% The argument should take the current |align_state| as input, and return the modified value (if any).
%    \begin{macrocode}
function getset(func)
    getset_func = func
    tex_runtoks(getset_run)
    getset_func = nil
    return getset_newval, getset_oldval
end
luaalignstate.getset = getset
end
%    \end{macrocode}
% \end{macro}
%
% \begin{macro}{
%     luaalignstate.get(),
%     luaalignstate.set(),
%     luaalignstate.add(),
% }
% Simple wrappers.
%    \begin{macrocode}
function luaalignstate.get() return select(1, getset()) end
do
local newval
local function retnewval()
    return newval
end
local function addvals(oldval)
    return oldval + newval
end
function luaalignstate.set(s)
    newval = s
    return getset(retnewval)
end
function luaalignstate.add(s)
    newval = s
    return getset(addvals)
end
end
%    \end{macrocode}
% \end{macro}
%
% Clean up and return.
%    \begin{macrocode}
node_direct_setfield(tmp_whatsit_ptr, "type", 100)  -- ("d"):byte()
node_direct_free(tmp_whatsit_ptr)
return luaalignstate
%</lua>
%    \end{macrocode}
%
%
% \subsection{\LaTeX{} package}
%
%    \begin{macrocode}
%<*package>
\ProvidesExplPackage
    {luaalignstate}
    {2026/10/01} {1}
    {Access the alignment counter in LuaTeX}
\sys_if_engine_luatex:F {
    \PackageError {luaalignstate} {LuaTeX~required} {}
    \file_input_stop:
}
%    \end{macrocode}
% In \LaTeX{}, we create the Lua functions the standard way, using |luatexbase|.
% To keep the Lua module \LaTeX{}-independent, we settle for a short block of inline Lua code.
%    \begin{macrocode}
\group_begin: \use:n {
    \char_set_catcode_space:n {32}
    \char_set_catcode_other:N \\
    \use:nnn
\group_end:
\lua_now:n}{{
    local tostring = tostring
    local tex_sprint = tex.sprint
    local token_scan_int = token.scan_int
    local token_set_lua = token.set_lua
    local token_set_macro = token.set_macro
    local lua_functions_table = lua.get_functions_table()
    local luatexbase_new_luafunction = luatexbase.new_luafunction
    luaalignstate = require"luaalignstate"
    local luaalignstate_getset = luaalignstate.getset
    local luaalignstate_add = luaalignstate.add
    _ENV = nil
    local function luacmd(name, f, ...)
        local id = luatexbase_new_luafunction(name)
        lua_functions_table[id] = f
        token_set_lua(name, id, "global", ...)
    end
%    \end{macrocode}
%
% \begin{macro}[EXP]{\luaalignstate_value:}
% The getter is straightforward.
%    \begin{macrocode}
    luacmd("luaalignstate_value:", function ()
        return tex_sprint(-2, tostring(luaalignstate_getset()))
    end)
%    \end{macrocode}
% \end{macro}
%
% \begin{macro}{
%     \@@_set:w,
%     \@@_add:w,
%     \luaalignstate_set:n,
%     \luaalignstate_add:n,
% }
% \begin{expl3codesyntax}
% \@@_set:w \<integer>
% \@@_add:w \<integer>
% \end{expl3codesyntax}
% For the setters, we make a luacall that uses primitive-like syntax, and a macro wrapper.
% The setter uses |getset()| rather than |set()|, to support accessing the old value during evaluation of the new value, by locally redefining \cs{luaalignstate_value:} (within the |\halign| group) to a token list containing the old value.
%    \begin{macrocode}
    local function set_scan(old_alignstate)
        token_set_macro("luaalignstate_value:", tostring(old_alignstate))
        return token_scan_int()
    end
    luacmd("@@_set:w", function ()
        return luaalignstate_getset(set_scan)
    end, "protected")
%    \end{macrocode}
% Note that $\cs{@@_set:w}$ doesn't scan its argument until \emph{inside} the |\halign|, whereas \cs{@@_add:w} scans its argument immediately, before initiating the |\halign|.
%    \begin{macrocode}
    luacmd("@@_add:w", function ()
        return luaalignstate_add(token_scan_int())
    end, "protected")
}}%
\cs_new_protected:Npn \luaalignstate_set:n #1 {
    \@@_set:w \tex_numexpr:D #1 \scan_stop:
}
\cs_new_protected:Npn \luaalignstate_add:n #1 {
    \@@_add:w \tex_numexpr:D #1 \scan_stop:
}
%    \end{macrocode}
% \end{macro}
%
% \begin{macro}{
%     \luaalignstate_get:N,
%     \luaalignstate_use:n,
%     \luaalignstate_use_inline:n,
% }
% For the in-context getters, we expand |\luaalignstate_value:| using brace tricks to balance the alignment counter.
%    \begin{macrocode}
\cs_new_protected:Npn \luaalignstate_get:N #1 {
    \tex_edef:D #1 { \if_false: } \fi: \luaalignstate_value: \if_false: { \fi: }
}
\cs_new:Npn \luaalignstate_use:n #1 {
    \tex_expanded:D {
        \tex_unexpanded:D {#1}
        { \if_false: }} \fi: \luaalignstate_value: \if_false: {{ \fi: }
    }
}
\cs_new:Npn \luaalignstate_use_inline:n #1 {
    \exp:w
    \tex_immediateassignment:D \cs_gset_nopar:Npn \@@_tmp:w ##1 {\exp_end: #1}
    \exp_after:wN \@@_tmp:w \tex_expanded:D {
        { \if_false: }} \fi: \luaalignstate_value: \if_false: {{ \fi: }
    }
}
%    \end{macrocode}
% \end{macro}
%
%    \begin{macrocode}
%</package>
%    \end{macrocode}
%
%
% \end{implementation}
